\documentclass[../../../../main.tex]{subfiles}
\begin{document}

\begin{flushright}
\footnotesize\textit{【自動文字起こし・要確認】}
\end{flushright}

{\bf [解]}
平面$\alpha$の点$X$として
\begin{align*}
\alpha:x+y+z=4
\end{align*}
だから，$\alpha$と$O$の間りの距離$L$は
\begin{align*}
L=\frac4{\sqrt3}=\frac{4\sqrt3}3<\sqrt6
\end{align*}
だから，$\alpha$と$S$は共有点を持つ．さらに，共有点の中心$A\left(\dfrac43,\dfrac43,\dfrac43\right)$，半径$\dfrac{\sqrt6}3$の円で，この円上の点は，$\alpha$に平行な単位ベクトル$\dfrac1{\sqrt2}\begin{pmatrix}0\\1\\-1\end{pmatrix}$，$\dfrac1{\sqrt6}\begin{pmatrix}2\\-1\\-1\end{pmatrix}$を用いて，
\begin{align*}
\overrightarrow{OX}=\frac43\begin{pmatrix}1\\1\\1\end{pmatrix}+\frac{\sqrt6}3\left[\frac1{\sqrt2}\cos\theta\begin{pmatrix}0\\1\\-1\end{pmatrix}+\frac1{\sqrt6}\sin\theta\begin{pmatrix}2\\-1\\-1\end{pmatrix}\right]\quad(0\le\theta<2\pi)
\end{align*}
とかけるから，$X(x,y,z)$として
\begin{align*}
x&=\frac43+\frac23\sin\theta\\
y&=\frac43-\frac{\sqrt3}3\cos\theta-\frac13\sin\theta\\
z&=\frac43+\frac{\sqrt3}3\cos\theta-\frac13\sin\theta
\end{align*}

以下$s=\sin\theta$，$c=\cos\theta$とすると
\begin{align*}
xyz&=\left(\frac43+\frac23s\right)\left\{\left(\frac43-\frac13s\right)^2-\frac13c^2\right\}\\
&=\frac23(2+s)\left\{\frac49s^2-\frac89s+\frac{13}9\right\}\\
&=\frac2{27}(4s^3-3s+26)
\end{align*}

$\{\ \}$の中身を$f(s)$とおく．
\begin{align*}
f'(s)=3(4s^2-1)
\end{align*}
から，下表を得る．

\begin{center}
\begin{tabular}{c|ccccccc}
$s$ & $-1$ & & $-\dfrac12$ & & $\dfrac12$ & & $1$ \\
\hline
$f'$ & & $+$ & $0$ & $-$ & $0$ & $+$ & \\
\hline
$f$ & & $\nearrow$ & & $\searrow$ & & $\nearrow$ & \\
\end{tabular}
\end{center}

ここで，
\begin{align*}
f(-1)=25，\quad f\left(\frac12\right)=25，\qquad f(1)=27，\quad f\left(-\frac12\right)=27
\end{align*}

から，もとめる積は
\begin{align*}
\frac{50}{27}\le xyz\le2．
\end{align*}

\bigskip
\noindent{\bf [解2]}
共有点上の点$(X,Y,Z)$は
\begin{align*}
X+Y+Z=4，\qquad X^2+Y^2+Z^2=6
\end{align*}
をみたす．$\alpha=XYZ$とおく．
\begin{align*}
YZ+ZX+XY=5
\end{align*}
から，$X,Y,Z$は$t$の3次式
\begin{align*}
t^3-4t^2+5t=\alpha
\end{align*}
の3実解．$f(t)=t^3-4t^2+5t$のグラフを考えると，$f'(t)=3t^2-8t+5=(3t-5)(t-1)$から，$t=1$で極大値$f(1)=2$，$t=5/3$で極小値$f(5/3)=\dfrac{50}{27}$をとるから，3実解をもつ条件は
\begin{align*}
\frac{50}{27}\le\alpha\le2
\end{align*}
となる．

\newpage

\end{document}
